\documentclass[11pt,letterpaper]{article}
\usepackage[utf8]{inputenc}
\usepackage[T1]{fontenc}
\usepackage[spanish,es-tabla,es-nodecimaldot,es-noquoting,es-noshorthands]{babel}
\usepackage[top=2.3cm,bottom=2.2cm,left=2.2cm,right=2.2cm]{geometry}
\usepackage{graphicx,booktabs,tabularx,array,xcolor,tikz,enumitem,fancyhdr,hyperref,amsmath,amssymb}

\definecolor{upvmorado}{HTML}{5B146B}
\definecolor{upvazul}{HTML}{17384D}
\definecolor{upvverde}{HTML}{008D79}
\definecolor{upvnaranja}{HTML}{E67817}
\definecolor{grisclaro}{HTML}{F8FAFC}
\definecolor{grisborde}{HTML}{CBD5E1}
\definecolor{gristexto}{HTML}{475569}

\hypersetup{
    colorlinks=true,
    urlcolor=upvverde,
    linkcolor=upvazul,
    pdftitle={Reporte Técnico - Proyecto final: Sistema inteligente para la generación de horarios de ITIID},
    pdfauthor={Equipo ITIID - Universidad Politécnica de Victoria}
}

\setlength{\parindent}{0pt}
\setlength{\parskip}{5pt}
\emergencystretch=1.8em
\setlist[itemize]{itemsep=2pt,topsep=3pt,leftmargin=*}
\setlist[enumerate]{itemsep=3pt,topsep=3pt,leftmargin=*}
\renewcommand{\arraystretch}{1.22}
\addto\captionsspanish{\renewcommand{\contentsname}{Índice General}}

\pagestyle{fancy}
\fancyhf{}
\fancyhead[L]{%
    \raisebox{-0.18\height}{\includegraphics[height=0.52cm]{img/logo_upv_morado.png}}\hspace{0.22cm}%
    {\footnotesize\bfseries\color{upvmorado} UPV $\cdot$ ITIID}\hspace{0.15cm}{\footnotesize\color{gristexto}| Reporte Técnico}%
}
\fancyhead[R]{\footnotesize\color{gristexto} Fundamentos de Inteligencia Artificial}
\fancyfoot[L]{\footnotesize\color{gristexto} Dr. José Fidencio López Luna}
\fancyfoot[R]{\footnotesize\bfseries\color{upvazul} Página \thepage}
\renewcommand{\headrulewidth}{0.5pt}
\renewcommand{\footrulewidth}{0.3pt}
\setlength{\headheight}{20pt}

\newcommand{\seccionUPV}[1]{%
    \section{\texorpdfstring{\textcolor{upvmorado}{#1}}{#1}}%
    \vspace{-0.22cm}{\color{upvnaranja}\rule{\textwidth}{1.1pt}}\vspace{0.12cm}%
}
\newcommand{\subseccionUPV}[1]{%
    \subsection{\texorpdfstring{\textcolor{upvazul}{#1}}{#1}}%
}

\begin{document}

\begin{titlepage}
    \centering
    \vspace*{0.1cm}
    \includegraphics[height=2.35cm]{img/logos_upv_bis.png}\\[0.35cm]
    {\color{upvmorado}\rule{\textwidth}{2pt}}\\[-0.22cm]
    {\color{upvnaranja}\rule{\textwidth}{0.9pt}}\\[0.35cm]

    {\Large\bfseries\color{upvazul} UNIVERSIDAD POLITÉCNICA DE VICTORIA\par}
    \vspace{0.12cm}
    {\normalsize\bfseries\color{gristexto} Ingeniería en Tecnologías de la Información e Innovación Digital (ITIID)\par}
    \vspace{0.35cm}

    \includegraphics[height=2.65cm]{img/logo_upv_morado.png}\\[0.35cm]

    {\LARGE\bfseries\color{upvmorado} Proyecto final: Sistema inteligente para la\\generación de horarios de ITIID\par}
    \vspace{0.28cm}
    {\large\bfseries\color{upvazul} Reporte Técnico de Arquitectura, Optimización Combinatoria\\e Integración de Inteligencia Artificial\par}
    \vspace{0.45cm}

    \begin{tikzpicture}
        \node[
            fill=grisclaro,
            draw=grisborde,
            rounded corners=4pt,
            inner xsep=14pt,
            inner ysep=9pt,
            text width=0.88\textwidth
        ] {%
            \begin{tabularx}{\linewidth}{@{}lX@{}}
                \textbf{\color{upvazul}Asignatura:} & Fundamentos de Inteligencia Artificial (7.\textsuperscript{o} Cuatrimestre) \\
                \textbf{\color{upvazul}Dr. encargado:} & \textbf{Dr. José Fidencio López Luna} \\
                \textbf{\color{upvazul}Institución:} & Universidad Politécnica de Victoria (UPV) --- Modelo BIS
            \end{tabularx}
        };
    \end{tikzpicture}

    \vspace{0.45cm}
    {\normalsize\bfseries\color{upvmorado} INTEGRANTES DEL EQUIPO\par}
    \vspace{0.18cm}

    \begin{tikzpicture}
        \node[
            fill=white,
            draw=upvmorado!35,
            rounded corners=4pt,
            inner xsep=16pt,
            inner ysep=10pt,
            text width=0.82\textwidth,
            align=center
        ] {%
            \begin{tabular}{@{}c@{}}
                Axel Librado Aguilera Alonso \\
                Germán David Zavala Castro \\
                Vasti Norely Ruiz Báez \\
                Yahir Alejandro Roque Martínez \\
                José Guadalupe Cornejo Alva \\
                David Alejandro Ibarra Castañeda \\
                Jesús Alberto Tovar Ramírez
            \end{tabular}
        };
    \end{tikzpicture}

    \vfill
    {\color{upvmorado}\rule{0.55\textwidth}{0.6pt}}\\[0.15cm]
    {\small\color{gristexto} Ciudad Victoria, Tamaulipas, México $\cdot$ Noviembre de 2026\par}
\end{titlepage}

\pagenumbering{roman}
\tableofcontents
\newpage
\pagenumbering{arabic}

\seccionUPV{Resumen Ejecutivo}

El presente reporte técnico documenta el diseño matemático, algorítmico y de ingeniería de software del \textbf{Sistema inteligente para la generación de horarios de ITIID} de la Universidad Politécnica de Victoria. El problema de calendarización universitaria (\textit{University Course Timetabling Problem}, UCTP) pertenece a la clase NP-completo y exige satisfacer simultáneamente restricciones duras institucionales (no traslape de docentes, grupos ni aulas; respeto estricto de disponibilidad, recesos y recursos compartidos; cobertura completa de sesiones semanales) mientras se minimizan penalizaciones blandas contrapuestas (horas muertas estudiantiles, huecos docentes, jornadas mayores a 6 horas y preferencias individuales).

La solución implementa una arquitectura híbrida gobernada por una regla estricta: \textbf{las asignaciones horarias son calculadas únicamente por motores formales y auditadas por un Validador Común determinista}, mientras que el \textbf{Modelo Grande de Lenguaje (LLM)} opera exclusivamente en tareas de traducción semántica con supervisión humana (\textit{human-in-the-loop}), explicación de conflictos y traducción de cambios imprevistos.

\seccionUPV{Formulación Matemática y Validador Común}

\subseccionUPV{Dominio Temporal y Entidades}
La semana académica consta de $D = 5$ días (lunes a viernes) con $S = 17$ franjas diarias de \texttt{07:00} a \texttt{20:45}, de las cuales tres corresponden a recesos fijos ($B = \{4, 8, 13\}$) y 14 son bloques lectivos ($H_{\text{lect}} = \{0,\dots,16\} \setminus B$), totalizando 70 franjas asignables por aula o grupo. El dimensionamiento automático de grupos por cuatrimestre $q$ responde a:
\[
G_q = \left\lceil \frac{\text{Alumnos}_q}{35} \right\rceil
\]

\subseccionUPV{Variables de Decisión y Restricciones Duras}
Sea $x_{c,s} \in \{0,1\}$ la variable binaria que vale $1$ si el curso $c \in C$ (con sesiones requeridas $w_c$, grupo $g_c$, docente $t_c$ y aula compatible $r_c$) se imparte en el bloque lectivo $s \in \mathcal{D}_c$:
\begin{align}
\sum_{s \in \mathcal{D}_c} x_{c,s} &= w_c && \forall c \in C \quad \text{(Cobertura semanal exacta)} \\
\sum_{c : g_c = g} x_{c,s} &\leq 1 && \forall g, \forall s \quad \text{(Sin choque de grupo)} \\
\sum_{c : t_c = t} x_{c,s} &\leq 1 && \forall t, \forall s \quad \text{(Sin choque de docente)} \\
\sum_{c : r_c = r} x_{c,s} &\leq 1 && \forall r, \forall s \quad \text{(Sin choque de aula compartida)}
\end{align}

\subseccionUPV{Función Objetivo Multicriterio y Escenarios A, B y C}
Para medir de manera lineal exacta las horas muertas diarias sin productos no lineales en CP-SAT, se introducen variables auxiliares booleanas monótonas $\text{st}_{e,d,p}$ (\textit{jornada iniciada}) y $\text{rm}_{e,d,p}$ (\textit{clases restantes}) sobre la secuencia compacta $p \in \{0,\dots,13\}$ de cada entidad $e$ (grupo o docente), de modo que existe hueco en $p$ si y solo si:
\[
\text{gap}_{e,d,p} \geq \text{st}_{e,d,p} + \text{rm}_{e,d,p} - 1 - a_{e,d,p}
\]
La función de penalización blanda pondera seis métricas bajo tres escenarios transparentes elegibles por coordinación académica:
\begin{center}
\small
\begin{tabular}{@{}lccc@{}}
\toprule
\textbf{Criterio Blando} & \textbf{Esc. A (Alumnos)} & \textbf{Esc. B (Docentes)} & \textbf{Esc. C (Equilibrio)} \\
\midrule
Horas muertas de alumnos & 10 & 2 & 6 \\
Jornadas excesivas ($>6$ h/día) & 8 & 3 & 5 \\
Concentración diaria de materia & 6 & 3 & 5 \\
Huecos entre clases docentes & 2 & 10 & 6 \\
Desbalance semanal docente & 2 & 8 & 5 \\
Preferencias docentes no cumplidas & 3 & 10 & 6 \\
\bottomrule
\end{tabular}
\end{center}

\seccionUPV{Motores de Resolución e Inteligencia Artificial}

\begin{enumerate}
    \item \textbf{R2 --- Google OR-Tools CP-SAT (\texttt{backend\_ia/solvers.py}):} Motor principal de programación por restricciones con satisfacción booleana, arranque cálido (\texttt{AddHint}), diagnóstico relajado automático ante instancias sobre-restringidas y soporte de \textbf{reprogramación estable (I4)} penalizando desviaciones respecto a un horario publicado ($x_{c,s} \neq x^{\text{orig}}_{c,s}$).
    \item \textbf{R3 --- Algoritmo Genético Multicorrida:} Representación directa por vector de sesiones, selección por torneo, cruce uniforme por curso, mutación guiada hacia dominios factibles y operador de reparación de colisiones, evaluado en 10 corridas independientes con reporte de media y desviación estándar.
    \item \textbf{R4 --- Búsqueda Local (MRV + Hill Climbing):} Construcción inicial priorizando cursos de menor dominio disponible (\textit{Minimum Remaining Values}) seguida de ascenso de colinas con múltiples reinicios aleatorios.
    \item \textbf{Motor Clásico C++ (\texttt{backend/solver.cpp}):} Búsqueda exacta por \textit{backtracking} con exploración horizontal por franja horaria y poda temprana de capacidad, disponibilidad y tope diario $\lceil w_c / 5 \rceil$.
    \item \textbf{R5 y R6 --- Servicio LLM y Privacidad (\texttt{backend\_ia/llm\_service.py}):} Integración con \textbf{Ollama Local} (\texttt{qwen2.5:3b}), APIs compatibles y respaldo determinista offline. Incluye anonimizador reversible (\texttt{Docente\_01}, \texttt{Aula\_01}), evaluación F1 contra dataset anotado (P2), explicador ejecutivo (P4) y auditoría empírica de generación directa con LLM (P1).
\end{enumerate}

\seccionUPV{Resultados Experimentales y Conclusiones}

En la comparativa estandarizada sobre la instancia cerrada \textbf{I2 (ITIID 7.\textsuperscript{o} Cuatrimestre, 7 materias, 35 sesiones semanales, Escenario C)}:
\begin{itemize}
    \item \textbf{OR-Tools CP-SAT (R2)} alcanza el óptimo global con \textbf{0 violaciones duras}, \textbf{0 horas muertas estudiantiles}, \textbf{0 huecos docentes}, \textbf{100\% de preferencias cumplidas} y penalización blanda mínima de \textbf{25 puntos}, superando ampliamente al horario manual de referencia (penalización 73, con 5 horas muertas).
    \item \textbf{Búsqueda Local (R4)} y \textbf{Algoritmo Genético (R3)} obtienen horarios 100\% factibles (\textbf{0 violaciones duras}) en menos de $2.3$ segundos con penalizaciones de 36 y 37 puntos respectivamente.
    \item En el escenario de imprevisto docente \textbf{I4} (bloqueo completo de miércoles), la reprogramación con CP-SAT repara la factibilidad moviendo únicamente \textbf{4 de 35 sesiones (88.6\% de estabilidad)}.
    \item La \textbf{Prueba P1} confirma empíricamente que un LLM directo sin motor de restricciones incurre entre 15 y 21 violaciones duras (choques y asignaciones en receso), validando la decisión arquitectónica de combinar optimización exacta con lenguaje natural supervisado.
\end{itemize}

\seccionUPV{Referencias Bibliográficas}
\begin{enumerate}
    \item Burke, E. K., \& Petrovic, S. (2002). Recent research directions in automated timetabling. \textit{European Journal of Operational Research}, 140(2), 266--280.
    \item Perron, L., \& Didier, F. (2024). \textit{CP-SAT Solver --- Google OR-Tools Documentation}. Google Optimization Tools.
    \item Bonutti, A., De Cesco, F., Di Gaspero, L., \& Schaerf, A. (2012). Benchmarking curriculum-based course timetabling: formulations, data formats, instances, validation, visualization, and results. \textit{Annals of Operations Research}, 194(1), 59--70.
    \item Cormen, T. H., Leiserson, C. E., Rivest, R. L., \& Stein, C. (2009). \textit{Introduction to Algorithms} (3rd ed.). MIT Press.
    \item Garey, M. R., \& Johnson, D. S. (1979). \textit{Computers and Intractability: A Guide to the Theory of NP-Completeness}. W. H. Freeman.
\end{enumerate}

\end{document}
